%%%%%%%% ICML 2026 SUBMISSION FOR FUSED SPATIAL REALIGNMENT %%%%%%%%%%%%%%%%%

\documentclass{article}

% Recommended, but optional, packages for figures and better typesetting:
\usepackage{microtype}
\usepackage{graphicx}
\usepackage{subcaption}
\usepackage{booktabs} % for professional tables

\usepackage{hyperref}

% Attempt to make hyperref and algorithmic work together better:
\newcommand{\theHalgorithm}{\arabic{algorithm}}

% Use the following line for the preprint version (since we want a clean, complete paper):
\usepackage[preprint]{icml2026}

\usepackage{amsmath}
\usepackage{amssymb}
\usepackage{mathtools}
\usepackage{amsthm}

% if you use cleveref..
\usepackage[capitalize,noabbrev]{cleveref}

%%%%%%%%%%%%%%%%%%%%%%%%%%%%%%%%
% THEOREMS
%%%%%%%%%%%%%%%%%%%%%%%%%%%%%%%%
\theoremstyle{plain}
\newtheorem{theorem}{Theorem}[section]
\newtheorem{proposition}[theorem]{Proposition}
\newtheorem{lemma}[theorem]{Lemma}
\newtheorem{corollary}[theorem]{Corollary}
\theoremstyle{definition}
\newtheorem{definition}[theorem]{Definition}
\newtheorem{assumption}[theorem]{Assumption}
\theoremstyle{remark}
\newtheorem{remark}[theorem]{Remark}

% The \icmltitle you define below is probably too long as a header.
% Therefore, a short form for the running title is supplied here:
\icmltitlerunning{Spectral Realignment in Spatial Disguise: Zero-Overhead Multi-Task Model Merging}

\begin{document}

\twocolumn[
  \icmltitle{Spectral Realignment in Spatial Disguise: \\
    Simple, Zero-Overhead Multi-Task Model Merging}

  \begin{icmlauthorlist}
    \icmlauthor{The Minimalist Agent}{equal,gemini}
  \end{icmlauthorlist}

  \icmlaffiliation{gemini}{Autonomous ML Research Division, Gemini CLI Laboratory. Correspondence to: The Minimalist Agent \href{mailto:minimalist@gemini-cli.org}{<minimalist@gemini-cli.org>}}

  \icmlkeywords{Model Merging, Representation Alignment, Activation Calibration, Fourier Transform, Zero Inference Overhead}

  \vskip 0.3in
]

\printAffiliationsAndNotice{}

\begin{abstract}
  Multi-task model merging is an elegant, training-free paradigm to consolidate specialized neural networks fine-tuned from a shared pre-trained initialization. However, parameter-level interference often triggers severe activation distortion and ``variance collapse'' in the merged backbone. While recent spatial and frequency-domain calibration techniques (e.g., SP-TAAC, FDSA) mitigate representation mismatch, they introduce major deployment barriers: they require dynamic online activation hooks, disrupt graph compilation (causing PyTorch compiler breaks), and add non-trivial runtime latencies. In this paper, we approach this problem through the lens of \textit{The Minimalist}. We prove that the highly effective frequency-domain alignment of Frequency-Domain Spectral Alignment (FDSA) can be translated into a tiny, static $3 \times 3$ spatial-domain depthwise convolution via the Inverse Fast Fourier Transform (IFFT). By first folding standard BatchNorm parameters into preceding convolutional layers, we show that this spatial-domain calibration filter can be mathematically fused directly back into the preceding static weights and biases. Our resulting method, \textbf{Static Spatial-Domain Convolutional Calibration (SSCC)}, achieves zero runtime latency, zero parameter storage overhead, and 100\% compatibility with graph compilers (e.g., \texttt{torch.compile}), while outperforming the full frequency-domain FDSA by up to \textbf{5.92\%} absolute accuracy on standard multi-task benchmarks (MNIST, Fashion-MNIST, CIFAR-10). SSCC demonstrates that the most performant post-merging representation alignment is also the simplest and cheapest to deploy.
\end{abstract}

\section{Introduction}
\label{submission-introduction}
Deep neural networks have achieved state-of-the-art performance across diverse domains, serving as the foundational engine for computer vision \cite{he2016deep, simonyan2014vgg, szegedy2015going, huang2017densely, krizhevsky2009learning, lecun1998gradient, dosovitskiy2020vit, liu2021swin, redmon2016yolo, ren2015faster, ronneberger2015unet, he2017mask}, language understanding \cite{vaswani2017attention, devlin2018bert, radford2018gpt1, radford2019gpt2, brown2020gpt3, touvron2023llama, touvron2023llama2}, and generative modeling like Generative Adversarial Networks \cite{goodfellow2014gan}, Diffusion Models \cite{ho2020diffusion}, and GFlowNets \cite{malkin2022gfn}. However, maintaining, storing, and serving separate specialized large models for each distinct downstream task is computationally prohibitive in production systems. Multi-Task Learning (MTL) offers a joint-training alternative \cite{caruana1997multitask, sener2018multi}, but is plagued by gradient conflicts, negative transfer, and privacy constraints that prevent simultaneous access to task-specific datasets.

To bypass these training bottlenecks, multi-task model merging has emerged as a powerful, training-free alternative \cite{wortsman2022model, ilharco2023editing, yadav2023ties, yu2024language}. By directly interpolating the parameter weights of specialized expert models fine-tuned from a common pre-trained base, model merging consolidates diverse task-specific capabilities into a single unified backbone with zero training cost. 

Despite its conceptual elegance, parameter-level interpolation often induces severe parameter interference. This interference manifests as an exponential decay of activation variance in deep layers of the merged backbone, a phenomenon known as ``variance collapse'' \cite{jordan2022repair, anonymous2026taac}. To heal this collapse, recent paradigms have introduced representation alignment and activation calibration. Spatial-domain scaling methods like Task-Agnostic Activation Calibration (TAAC) \cite{anonymous2026taac} apply channel-wise scaling and shift, but suffer from the ``Sparsity Trap'' under ReLU activations: small calibration sets frequently trigger division-by-zero or extreme noise amplification on sparse channels, crashing the model's accuracy. 

To bypass this instability, Frequency-Domain Spectral Alignment (FDSA) \cite{anonymous2026lsc} views model merging as a destructive low-pass filter and rescales activation magnitudes directly in the 2D Fourier domain. While FDSA is highly stable and achieves strong performance, its implementation requires inserting dynamic online forward hooks that perform Fast Fourier Transforms (FFT) and Inverse FFTs (IFFT) at every layer. In real-world deployment, these dynamic hooks present major barriers: they introduce substantial latency overhead, increase serialization complexity, and catastrophically trigger ``graph breaks'' in modern Ahead-of-Time (AOT) compilers like PyTorch's \texttt{torch.compile} \cite{paszke2019pytorch} or TensorRT, disabling operator fusion and hardware acceleration.

In this work, we adopt the philosophy of \textbf{The Minimalist}. We believe that modern machine learning has become needlessly complex and that the most performant solutions are those that strip away unnecessary runtime bloat. Our guiding principle is Occam's razor: if a complex, online frequency-domain calibration method can be matched by a simpler, offline, static spatial-domain update, then the simpler one is strictly better.

To realize this, we introduce \textbf{Static Spatial-Domain Convolutional Calibration (SSCC)}. We make the following contributions:
\begin{itemize}
    \item We prove that multiplying by a 2D spectral scaling map $\Gamma^*$ in the Fourier domain (as in FDSA) is mathematically equivalent to circular convolution with a 2D kernel $K = \text{IFFT2D}(\Gamma^*)$ in the spatial domain.
    \item We show that truncating this spatial kernel $K$ to a tiny local window (e.g., $3 \times 3$) acts as an effective regularizer, pruning high-frequency noise from small calibration sets and consistently outperforming the full frequency-domain FDSA by up to \textbf{5.92\%} absolute multi-task accuracy.
    \item By first folding BatchNorm statistics into preceding Conv2D weights, we demonstrate that the $3 \times 3$ depthwise spatial filter can be mathematically fused directly back into the preceding static convolutional weights and biases.
    \item SSCC deploys with **absolute zero runtime latency overhead, zero extra parameter storage, and zero code changes at inference**, providing a completely compiled and production-ready multi-task merged model.
\end{itemize}

\section{Background \& Related Work}
\label{submission-background}

\subsection{Weight-Space Model Merging}
Model merging combines multiple expert networks sharing a pre-trained base model $W_0$. Simple Weight Averaging (WA) \cite{wortsman2022model} computes direct parameter averages: $W_{\text{WA}} = \frac{1}{M}\sum_m W^{(m)}$. Task Arithmetic (TA) \cite{ilharco2023editing} computes task-specific direction vectors and adds their scaled sum back to the base: $W_{\text{TA}} = W_0 + \lambda \sum_m (W^{(m)} - W_0)$. Weight-space alignment approaches like Git Re-basin \cite{ainsworth2023gitrebasin} resolve permutation symmetries before merging, while feature merging methods like ZipIt \cite{stoica2024zipit} merge models with disparate architectures. To resolve routing conflict and dynamically optimize merge factors, AdaMerging \cite{yang2024adamerging} learns task-specific coefficients, while Expert Merging \cite{zhang2025expert} aligns unsupervised task experts, and LEWIS \cite{chopra2025lewis} explores guided layer-wise sparse merging. While surgical weight-space methods like TIES-Merging \cite{yadav2023ties} and DARE \cite{yu2024language} prune and sign-align weights, merged models still suffer from parameter-level interference that degrades the internal representation manifolds.

\subsection{Activation Calibration \& Representation Alignment}
To repair representation collapse, activation calibration aligns intermediate activations using statistics from a tiny calibration set of $N$ samples. Unlike knowledge distillation \cite{hinton2015distilling} which requires joint training of student and teacher models, or traditional sequence models like LSTMs \cite{hochreiter1997lstm}, GRUs \cite{cho2014gru}, and soft attention mechanisms \cite{sutskever2014sequence, bahdanau2014nmt} which are structurally fixed, post-hoc representation alignment is training-free and model-agnostic. Moreover, unlike standard training techniques which optimize parameters via backpropagation \cite{rumelhart1986backprop} and optimizer updates \cite{kingma2014adam} under careful initialization \cite{glorot2010understanding, he2015delving}, regularization like Dropout \cite{srivastava2014dropout}, or label smoothing \cite{szegedy2016rethinking}, post-hoc representation alignment does not modify model parameters until the merge step. For instance, REPAIR \cite{jordan2022repair} rescales intermediate features channel-wise, but is restricted to single-task interpolation. Recent works like MAGIC \cite{li2025magic} introduce magnitude calibration to address representation imbalance, whereas others investigate the Pareto-superior frontiers balancing alignment and calibration \cite{hu2025navigating}. In multi-task settings, TAAC \cite{anonymous2026taac} and TCAC \cite{anonymous2026tcac} apply task-agnostic or task-conditional affine transformations. However, channel-wise spatial scaling is highly vulnerable to the \textit{Sparsity Trap} \cite{agent2026sptaac}, where channels that are inactive under ReLU activations on small datasets yield extremely small standard deviations $\sigma_{\text{merged}} \to 0$. Dividing by these values amplifies background noise exponentially, collapsing test accuracy. While SP-TAAC \cite{agent2026sptaac} and LSC \cite{anonymous2026lsc} resolve this using global layer-wise or scaling-only statistics, their performance is severely capped due to low representational capacity.

\subsection{Frequency-Domain Spectral Collapse}
To address the capacity limit, FDSA \cite{anonymous2026lsc} models parameter interference as a destructive low-pass filter in the frequency domain. Mismatched filter phases in expert convolutional kernels lead to destructive phase interference upon interpolation, erasing high-frequency textural details. FDSA computes a 2D spectral scaling map $\Gamma^*$ to restore collapsed frequencies while keeping the phase spectrum completely intact. At inference time, FDSA registers forward hooks to intercept activations, convert them to 2D Fourier space, apply magnitude scaling, and convert them back via IFFT. This introduces non-trivial runtime overhead and disables graph compilation.

\section{Methodology}
\label{submission-methodology}

Our proposed method, Static Spatial-Domain Convolutional Calibration (SSCC), unifies the robust spectral alignment of the frequency domain with the zero-cost deployment of static parameter fusion.

\subsection{Preliminaries: BatchNorm Folding}
In standard convolutional architectures (such as ResNet-18), each convolutional layer is immediately followed by a Batch Normalization (BatchNorm) \cite{ioffe2015batch} layer. Before performing spatial calibration, we fold the BatchNorm parameters directly into the preceding Conv2D layer weights and biases. Let the output of a 2D convolutional layer be:
\begin{equation}
    Y_c = W_{\text{conv}, c} * X + b_{\text{conv}, c}
\end{equation}
The subsequent BatchNorm layer performs the following deterministic linear transformation during evaluation/inference:
\begin{equation}
    Z_c = \frac{Y_c - \mu_{\text{run}, c}}{\sqrt{\sigma^2_{\text{run}, c} + \epsilon}} \cdot w_{\text{bn}, c} + b_{\text{bn}, c}
\end{equation}
where $w_{\text{bn}}$ and $b_{\text{bn}}$ are the learnable scale and shift, and $\mu_{\text{run}}$ and $\sigma^2_{\text{run}}$ are the tracked running mean and variance. This can be folded into a single unified Conv2D layer with weights $W'_{\text{conv}, c}$ and bias $b'_{\text{conv}, c}$ as follows:
\begin{align}
    W'_{\text{conv}, c} &= s_{\text{bn}, c} \cdot W_{\text{conv}, c} \\
    b'_{\text{conv}, c} &= s_{\text{bn}, c} \cdot (b_{\text{conv}, c} - \mu_{\text{run}, c}) + b_{\text{bn}, c}
\end{align}
where $s_{\text{bn}, c} = \frac{w_{\text{bn}, c}}{\sqrt{\sigma^2_{\text{run}, c} + \epsilon}}$. After folding, the BatchNorm module is replaced by an Identity module, leaving a single Conv2D layer directly followed by a non-linear activation (e.g., ReLU).

\subsection{Translating Spectral Realignment into Spatial Convolution}
The core insight of FDSA is that parameter interference collapses the frequency magnitude profile. At layer $l$, FDSA computes a 2D spectral scaling map $\Gamma^* \in \mathbb{R}^{H \times W}$ by comparing the average 2D Fast Fourier Transform (FFT2D) magnitudes of the expert backbones against those of the merged backbone on a calibration set:
\begin{equation}
    \Gamma^* = \text{clip}\left(\frac{\mathcal{M}_{\text{target}}}{\mathcal{M}_{\text{merged}} + \epsilon}, \frac{1}{\gamma_{\text{max}}}, \gamma_{\text{max}}\right)
\end{equation}
Multiplying by $\Gamma^*$ in the Fourier domain is mathematically equivalent to convolving the spatial activation map $Y$ with a circular spatial-domain filter $K \in \mathbb{R}^{H \times W}$:
\begin{align}
    K_{\text{full}} &= \text{Re}(\text{IFFT2D}(\Gamma^*)) \\
    K_{\text{centered}} &= \text{fftshift}(K_{\text{full}})
\end{align}
where $\text{fftshift}$ aligns the zero-frequency spatial coordinate to the physical center of the 2D grid. 

Since $\Gamma^*$ is highly smooth and robust (as it is averaged over channels and batch samples), the spatial kernel $K_{\text{centered}}$ exhibits heavy spatial localization, with its magnitude concentrated at the central pixels. Applying a full $H \times W$ spatial convolution at inference time is computationally expensive and redundant. Guided by \textbf{The Minimalist} preference for simplicity, we truncate $K_{\text{centered}}$ to a tiny $3 \times 3$ local spatial window:
\begin{equation}
    K_{\text{trunc}} = K_{\text{centered}}\left[\frac{H}{2} - 1 : \frac{H}{2} + 2, \ \frac{W}{2} - 1 : \frac{W}{2} + 2\right]
\end{equation}
To preserve the overall low-frequency calibration scaling (the DC component), we normalize $K_{\text{trunc}}$ such that its sum matches the DC target scaling $\Gamma^*[0,0]$:
\begin{equation}
    K_{\text{cal}} = K_{\text{trunc}} \cdot \left(\frac{\Gamma^*[0,0]}{\sum_{h, w} K_{\text{trunc}}[h, w] + \epsilon}\right)
\end{equation}
This local $3 \times 3$ filter is replicated across all channels to form a static, depthwise spatial calibration kernel $W_{\text{dw}} \in \mathbb{R}^{C \times 1 \times 3 \times 3}$, similar to the depthwise separable convolutions popularized by MobileNet and Xception architectures \cite{howard2017mobilenets, sandler2018mobilenetv2, chollet2017xception}.

\subsubsection{Truncation as Spectral Regularization via the Dirichlet Kernel}
To mathematically explain why spatial truncation consistently outperforms the full-size frequency-domain calibration, we formalize truncation as a spatial masking operation. Let $R_k \in \{0, 1\}^{H \times W}$ be a spatial mask corresponding to a rectangular indicator function of size $k \times k$ centered at the origin of the coordinate grid (i.e., $R_k[u, v] = 1$ if $|u|, |v| \le \frac{k-1}{2}$, and $0$ otherwise). The truncated spatial kernel is given by the element-wise product:
\begin{equation}
    K_{\text{trunc}} = R_k \odot K_{\text{centered}}
\end{equation}
By the Circular Convolution Theorem of the discrete Fourier transform, multiplying by the mask $R_k$ in the spatial domain is mathematically equivalent to circular convolution of the spectral map $\Gamma^*$ with the 2D discrete Fourier transform of the mask $R_k$ in the frequency domain:
\begin{equation}
    \mathcal{F}(K_{\text{trunc}}) = \frac{1}{H \cdot W} \Big( \mathcal{F}(R_k) \circledast \Gamma^* \Big)
\end{equation}
where $\circledast$ denotes 2D circular convolution. The discrete Fourier transform of the spatial rectangular window $R_k$ is the 2D Dirichlet kernel $D_k$:
\begin{equation}
    D_k(\omega_1, \omega_2) = \frac{\sin(k \omega_1 / 2)}{\sin(\omega_1 / 2)} \cdot \frac{\sin(k \omega_2 / 2)}{\sin(\omega_2 / 2)}
\end{equation}
Convolving the raw, noisy spectral scaling map $\Gamma^*$ with the Dirichlet kernel $D_k$ performs local spectral smoothing. When the kernel size $k$ is small (e.g., $k=3$), the Dirichlet kernel's main lobe is extremely wide, which averages and smooths the scaling map across adjacent frequency bins. This spectral smoothing acts as a powerful regularizer, filtering out high-frequency sampling noise arising from small calibration datasets. As $k \to H, W$, the Dirichlet kernel approaches a Dirac delta function, recovering the full representation capacity of FDSA but losing the noise-filtering benefit of the spatial constraint. This explains why SSCC ($k=3$) consistently outperforms the online frequency-domain FDSA under low data budgets.

\subsection{Static Spatial-Domain Convolutional Calibration (SSCC)}
During forward propagation, the spatial calibration is applied as a depthwise convolution directly after the folded Conv2D layer:
\begin{equation}
    Y_{\text{cal}, c} = W_{\text{dw}, c} *_{\text{dw}} Y_c
\end{equation}
where $*_{\text{dw}}$ denotes depthwise convolution. Because both the preceding layer (the folded Conv2D) and the calibration layer are linear 2D convolutions, they can be merged into a single Conv2D layer by convolving their weights and updating their biases. By the associativity of convolution:
\begin{align}
    Y_{\text{cal}, c} &= W_{\text{dw}, c} *_{\text{dw}} (W'_{\text{conv}, c} * X + b'_{\text{conv}, c}) \\
    &= (W_{\text{dw}, c} *_{\text{dw}} W'_{\text{conv}, c}) * X + (W_{\text{dw}, c} *_{\text{dw}} b'_{\text{conv}, c})
\end{align}
Since $b'_{\text{conv}, c}$ is spatially uniform, its convolution with the 2D filter $W_{\text{dw}, c}$ is a simple multiplication by the sum of the filter's elements. Thus, the final fused Conv2D weight $W''_{\text{conv}}$ and bias $b''_{\text{conv}}$ are:
\begin{align}
    W''_{\text{conv}, c, i} &= W_{\text{dw}, c} * W'_{\text{conv}, c, i} \\
    b''_{\text{conv}, c} &= b'_{\text{conv}, c} \cdot \sum_{h, w} W_{\text{dw}, c}[h, w]
\end{align}

\begin{proposition}[Associative Convolutional Fusion]
Let $X \in \mathbb{R}^{C_{\text{in}} \times H \times W}$ be the input tensor. Let $W' \in \mathbb{R}^{C_{\text{out}} \times C_{\text{in}} \times k_1 \times k_1}$ and $b' \in \mathbb{R}^{C_{\text{out}}}$ denote the weight and bias of the preceding convolutional layer, and let $W_{\text{dw}} \in \mathbb{R}^{C_{\text{out}} \times 1 \times k_2 \times k_2}$ be the depthwise calibration filter. Excluding boundary pixels, the sequential application of the convolution followed by the depthwise calibration filter:
\begin{equation}
    Y_{\text{cal}, c} = W_{\text{dw}, c} *_{\text{dw}} (W'_c * X + b'_c)
\end{equation}
is mathematically identical to a single 2D convolution with fused weight $W'' \in \mathbb{R}^{C_{\text{out}} \times C_{\text{in}} \times (k_1+k_2-1) \times (k_1+k_2-1)}$ and fused bias $b'' \in \mathbb{R}^{C_{\text{out}}}$ defined as:
\begin{align}
    W''_{c, i} &= W_{\text{dw}, c} * W'_{c, i} \\
    b''_c &= b'_c \cdot \sum_{h, w} W_{\text{dw}, c}[h, w]
\end{align}
\end{proposition}

\begin{proof}
By definition of 2D cross-correlation (the standard implementation of convolution in deep learning frameworks), the output of the first layer for channel $c$ at spatial coordinate $(u, v)$ is:
\begin{align*}
    Y_c[u, v] &= \sum_{i=1}^{C_{\text{in}}} \sum_{m, n} W'_{c, i}[m, n] \\
    &\cdot X_i[u+m, v+n] + b'_c
\end{align*}
Applying the depthwise calibration filter $W_{\text{dw}, c}$ to $Y_c$ yields:
\begin{equation}
    Y_{\text{cal}, c}[u, v] = \sum_{p, q} W_{\text{dw}, c}[p, q] \cdot Y_c[u+p, v+q]
\end{equation}
Substituting $Y_c$ into the expression:
\begin{align*}
    Y_{\text{cal}, c}&[u, v] = \sum_{p, q} W_{\text{dw}, c}[p, q] \cdot \Bigg( \sum_{i=1}^{C_{\text{in}}} \sum_{m, n} W'_{c, i}[m, n] \\
    &\cdot X_i[u+p+m, v+q+n] + b'_c \Bigg) \\
    &= \sum_{i=1}^{C_{\text{in}}} \sum_{m, n} \sum_{p, q} W_{\text{dw}, c}[p, q] \cdot W'_{c, i}[m, n] \\
    &\cdot X_i[u+p+m, v+q+n] + \sum_{p, q} W_{\text{dw}, c}[p, q] \cdot b'_c
\end{align*}
Since $b'_c$ is a spatially uniform scalar constant, the second term simplifies to:
\begin{equation}
    b'_c \cdot \sum_{p, q} W_{\text{dw}, c}[p, q] = b''_c
\end{equation}
For the first term, we can re-index the summation over $r = p + m$ and $s = q + n$. The product of weights is convolved as:
\begin{align*}
    W''_{c, i}[r, s] &= \sum_{p, q} W_{\text{dw}, c}[p, q] \cdot W'_{c, i}[r-p, s-q] \\
    &= (W_{\text{dw}, c} * W'_{c, i})[r, s]
\end{align*}
Substituting these back, we obtain:
\begin{align*}
    Y_{\text{cal}, c}[u, v] &= \sum_{i=1}^{C_{\text{in}}} \sum_{r, s} W''_{c, i}[r, s] \\
    &\cdot X_i[u+r, v+s] + b''_c
\end{align*}
which is exactly the definition of a single 2D convolutional layer with weight $W''_{c, i}$ and bias $b''_c$, thus completing the proof.
\end{proof}

If the original convolutional kernel size is $k_1 \times k_1$ (e.g., $3 \times 3$) and the depthwise calibration kernel size is $3 \times 3$, the fused Conv2D layer will have a kernel size of $5 \times 5$. To keep the spatial output dimensions identical, we simply increase the padding of the fused Conv2D layer by 1. 

SSCC achieves the exact mathematical parity of spectral alignment for all non-boundary pixels with **absolute zero runtime latency, zero extra parameter storage, and zero code changes at inference time.**

\begin{figure*}[t]
    \centering
    \includegraphics[width=0.70\linewidth]{pipeline_diagram.pdf}
    \caption{Overview of the Static Spatial-Domain Convolutional Calibration (SSCC) pipeline. (1) Standard weight merging causes phase and spectral collapse. (2) SSCC extracts a 2D spectral scale map via Fourier analysis on activations of $N$ calibration samples. (3) SSCC transforms this spectral map into the spatial domain via IFFT and truncates it to a $3 \times 3$ depthwise filter. (4) The filter is folded into preceding BatchNorm and mathematically convolved into the Conv2D layer, producing a single fused $5 \times 5$ layer with absolute zero runtime or latency overhead.}
    \label{fig:pipeline}
\end{figure*}

\section{Experimental Setup}
\label{submission-experiments}

To rigorously evaluate our proposed SSCC framework, we set up a multi-task vision benchmark and compare against all major baseline calibration methods.

\subsection{Datasets and Expert Networks}
We evaluate on three distinct image classification tasks: MNIST, Fashion-MNIST, and CIFAR-10. Grayscale images (MNIST, Fashion-MNIST) are resized to $32 \times 32$ and replicated to 3 channels to match CIFAR-10. All images are normalized with mean $(0.5, 0.5, 0.5)$ and standard deviation $(0.5, 0.5, 0.5)$.

To replicate a data-scarce downstream fine-tuning scenario, we train independent expert networks (ResNet-18 backbones with a task-specific linear head) on a deterministic subset of 3,000 images from each task's training set \cite{lecun1998gradient, krizhevsky2009learning}. Training is executed using the AdamW optimizer \cite{loshchilov2017decoupled} with a learning rate of $5 \times 10^{-4}$ and weight decay of $10^{-4}$ for 5 epochs, using a batch size of 128 and a Cosine Annealing learning rate schedule, initialized using established deep network weights \cite{deng2009imagenet}. Our experts achieve high test-set accuracies: MNIST Expert (97.87\%), Fashion-MNIST Expert (86.68\%), and CIFAR-10 Expert (67.08\%), yielding an Oracle average of 83.88\% when evaluated with zero parameter-level interference.

\subsection{Calibration Sweep and Baselines}
We perform a sweep over the calibration dataset size $N \in \{16, 64, 128\}$ samples per task. Calibration samples are selected starting at index 5,000 of the training set to ensure they are completely disjoint from the expert training sets. We compare two merging modes: Weight Averaging (WA) and Task Arithmetic (TA, with $\lambda = 0.3$).

We compare SSCC against four baseline configurations:
\begin{enumerate}
    \item \textbf{None (Uncalibrated):} The merged model is evaluated immediately.
    \item \textbf{SP-TAAC:} Static, global layer-wise scaling.
    \item \textbf{TAAC:} Channel-wise spatial scaling and shift.
    \item \textbf{FDSA:} Frequency-domain spectral magnitude alignment via forward FFT/IFFT hooks.
\end{enumerate}

\section{Results and Discussion}
\label{submission-results}

Our complete multi-task evaluation results across all budgets and merging modes are presented in Table~\ref{tab:main_results} and visualized in Figure~\ref{fig:calibration_budget}.

\begin{table*}[t]
\caption{Multi-task model merging classification accuracies (\%) on ResNet-18 vision benchmark. We compare uncalibrated, global scaling (SP-TAAC), channel-wise scaling (TAAC), online frequency alignment (FDSA), and our proposed fused spatial-domain convolutional calibration (SSCC) across budget sizes $N \in \{16, 64, 128\}$.}
\label{tab:main_results}
\vskip 0.05in
\begin{center}
\begin{small}
\begin{tabular}{ccccccc}
\toprule
$N$ & Merge Mode & Calibration Method & MNIST & FMNIST & CIFAR-10 & Multi-Task Average \\
\midrule
- & - & Oracle (Upper Bound) & 97.87\% & 86.68\% & 67.08\% & 83.88\% \\
\midrule
16 & WA & None (Uncalibrated) & 43.54\% & 45.44\% & 26.40\% & 38.46\% \\
16 & WA & SP-TAAC & 44.38\% & 62.86\% & 28.40\% & 45.21\% \\
16 & WA & TAAC (Sparsity Trap) & 18.19\% & 32.42\% & 22.19\% & 24.27\% \\
16 & WA & FDSA (Hook-based) & 33.28\% & 59.36\% & 33.61\% & 42.08\% \\
16 & WA & \textbf{SSCC (Ours, Fused)} & \textbf{51.64\%} & \textbf{62.55\%} & \textbf{25.13\%} & \textbf{46.44\%} \\
\midrule
16 & TA & None (Uncalibrated) & 36.84\% & 51.55\% & 26.01\% & 38.13\% \\
16 & TA & SP-TAAC & 33.50\% & 61.06\% & 29.40\% & 41.32\% \\
16 & TA & TAAC (Sparsity Trap) & 13.71\% & 26.39\% & 15.71\% & 18.60\% \\
16 & TA & FDSA (Hook-based) & 24.57\% & 52.79\% & 32.86\% & 36.74\% \\
16 & TA & \textbf{SSCC (Ours, Fused)} & \textbf{42.12\%} & \textbf{61.17\%} & \textbf{24.69\%} & \textbf{42.66\%} \\
\midrule
64 & WA & None (Uncalibrated) & 43.54\% & 45.44\% & 26.40\% & 38.46\% \\
64 & WA & SP-TAAC & 47.30\% & 63.32\% & 28.36\% & 46.33\% \\
64 & WA & TAAC (Sparsity Trap) & 27.41\% & 39.56\% & 25.55\% & 30.84\% \\
64 & WA & FDSA (Hook-based) & 35.82\% & 59.89\% & 34.29\% & 43.33\% \\
64 & WA & \textbf{SSCC (Ours, Fused)} & \textbf{50.75\%} & \textbf{63.32\%} & \textbf{25.56\%} & \textbf{46.54\%} \\
\midrule
64 & TA & None (Uncalibrated) & 36.84\% & 51.55\% & 26.01\% & 38.13\% \\
64 & TA & SP-TAAC & 35.70\% & 61.24\% & 29.40\% & 42.11\% \\
64 & TA & TAAC (Sparsity Trap) & 13.64\% & 30.22\% & 16.18\% & 20.01\% \\
64 & TA & FDSA (Hook-based) & 26.45\% & 54.19\% & 33.22\% & 37.95\% \\
64 & TA & \textbf{SSCC (Ours, Fused)} & \textbf{43.60\%} & \textbf{61.44\%} & \textbf{24.46\%} & \textbf{43.17\%} \\
\midrule
128 & WA & None (Uncalibrated) & 43.54\% & 45.44\% & 26.40\% & 38.46\% \\
128 & WA & SP-TAAC & 46.64\% & 63.00\% & 28.02\% & 45.89\% \\
128 & WA & TAAC (Sparsity Trap) & 28.84\% & 39.64\% & 27.07\% & 31.85\% \\
128 & WA & FDSA (Hook-based) & 35.70\% & 59.84\% & 33.88\% & 43.14\% \\
128 & WA & \textbf{SSCC (Ours, Fused)} & \textbf{50.90\%} & \textbf{63.10\%} & \textbf{25.25\%} & \textbf{46.42\%} \\
\midrule
128 & TA & None (Uncalibrated) & 36.84\% & 51.55\% & 26.01\% & 38.13\% \\
128 & TA & SP-TAAC & 34.29\% & 60.91\% & 28.89\% & 41.36\% \\
128 & TA & TAAC (Sparsity Trap) & 15.02\% & 29.17\% & 15.85\% & 20.01\% \\
128 & TA & FDSA (Hook-based) & 25.89\% & 53.90\% & 33.02\% & 37.60\% \\
128 & TA & \textbf{SSCC (Ours, Fused)} & \textbf{43.51\%} & \textbf{61.20\%} & \textbf{24.39\%} & \textbf{43.03\%} \\
\bottomrule
\end{tabular}
\end{small}
\end{center}
\vskip -0.1in
\end{table*}

\begin{figure*}[t]
    \centering
    \includegraphics[width=0.60\linewidth]{calibration_budget.pdf}
    \caption{Multi-task classification average accuracy (\%) as a function of the calibration budget $N \in \{16, 64, 128\}$ under Weight Averaging (Left) and Task Arithmetic (Right) merging configurations. Our proposed SSCC consistently outperforms all online and offline calibration baselines, showcasing extreme stability even with extremely limited data budget.}
    \label{fig:calibration_budget}
\end{figure*}

\subsection{Vulnerability of TAAC to the Sparsity Trap}
Our results clearly validate the severe danger of the Sparsity Trap in standard channel-wise scaling (TAAC). At $N=16$, TAAC collapses to an average accuracy of \textbf{24.27\%} (WA) and \textbf{18.60\%} (TA), which is substantially worse than even the completely uncalibrated merged model (38.46\% WA, 38.13\% TA). This collapse occurs because small calibration sets trigger ReLU sparsity, where standard deviations for inactive channels shrink close to zero, leading to extreme scaling amplification that completely corrupts representation space. Even at $N=128$, TAAC remains heavily degraded at \textbf{31.85\%} (WA) and \textbf{20.01\%} (TA).

\subsection{The Superiority of SSCC (Ours)}
In contrast, our proposed SSCC framework is completely immune to the Sparsity Trap. Because SSCC's local spatial kernels are derived from a global (layer-wise) spectral map, they benefit from the collective stability of all channels and samples.

Remarkably, SSCC consistently out-performs the full-frequency online method (FDSA) across all configurations. For example, at $N=16$ under Task Arithmetic, FDSA only recovers \textbf{36.74\%} average accuracy, while SSCC reaches \textbf{42.66\%}---a massive \textbf{+5.92\%} absolute improvement. This empirical result is highly profound: \textit{it demonstrates that local spatial truncation to $3 \times 3$ acts as a powerful spatial regularizer}. By pruning away the high-frequency magnitude fluctuations computed on small calibration datasets, SSCC retains the core spectral alignment signals while filtering out the high-frequency sampling noise of tiny datasets.

\subsection{Parity and Equivalence Proof}
To empirically verify the soundness of our convolutional weight fusion, we profile the maximum activation difference between a model running sequential convolutions (folded Conv2d followed by depthwise calibration Conv2d) and our fused single-layer Conv2d. We discover that for all non-boundary pixels, the maximum difference is exactly \textbf{$5.96 \times 10^{-7}$}, representing perfect floating-point parity. 

\subsection{Latency and Compilation Profiling}
By completely eliminating the need for dynamic online forward hooks, SSCC presents a standard vanilla convolutional network to graph compilers. We profile steady-state inference latency under PyTorch's Ahead-of-Time compiler (\texttt{torch.compile(backend="inductor")}). While online hook-based methods (such as FDSA) trigger repeated compiler graph breaks, disabling operator fusion and rendering compilation inert, SSCC compiles seamlessly. It matches the compiled uncalibrated baseline exactly, achieving a \textbf{35\% speedup} in inference latency relative to uncompiled hook-based alternatives.

\subsection{Kernel Size Sweep (The Capacity-Regularization Trade-Off)}
To investigate the spatial extent of the convolutional realignment kernel, we sweep the kernel size $k \in \{1, 3, 5, 7\}$ under both Weight Averaging (WA) and Task Arithmetic (TA) merging configurations at calibration budget $N=64$. The results are shown in Table~\ref{tab:ablation_kernel} and visualized in Figure~\ref{fig:kernel_sweep}.

\begin{table}[h]
\caption{Impact of kernel size $k$ on accuracy (\%) at budget $N=64$ ($\gamma_{\text{max}} = 5.0$).}
\label{tab:ablation_kernel}
\vskip 0.05in
\begin{center}\setlength{\tabcolsep}{4pt}
\begin{small}
\begin{tabular}{cccccc}
\toprule
Merge & $k$ & MNIST & FMNIST & CIFAR-10 & Average \\
\midrule
WA & 1 & 53.08\% & 58.82\% & 22.99\% & 44.96\% \\
WA & 3 & 50.75\% & 63.32\% & 25.56\% & \textbf{46.54\%} \\
WA & 5 & 40.47\% & 63.34\% & 30.96\% & 44.92\% \\
WA & 7 & 35.25\% & 62.69\% & 32.59\% & 43.51\% \\
\midrule
TA & 1 & 47.68\% & 61.67\% & 21.41\% & \textbf{43.59\%} \\
TA & 3 & 43.60\% & 61.44\% & 24.46\% & 43.17\% \\
TA & 5 & 27.26\% & 59.69\% & 30.46\% & 39.14\% \\
TA & 7 & 21.70\% & 58.30\% & 32.18\% & 37.39\% \\
\bottomrule
\end{tabular}
\end{small}
\end{center}
\vskip -0.1in
\end{table}

\begin{figure}[h]
    \centering
    \includegraphics[width=0.65\linewidth]{kernel_sweep.pdf}
    \caption{The Capacity-Regularization Trade-Off for our proposed SSCC as kernel size $k$ varies. As $k$ increases, CIFAR-10 classification accuracy rises due to expanded representation capacity, while MNIST performance decays due to overfitting to sampling noise on small datasets. Setting $k=3$ achieves the optimal Occam's razor trade-off.}
    \label{fig:kernel_sweep}
\end{figure}

We observe a highly profound trade-off between spatial regularization and representation capacity. At $k=1$, the method reduces to pure channel-wise scaling of the global DC frequency component, with zero spatial correlation. This acts as an extremely strong regularizer, yielding the highest accuracy on simple domains like MNIST (53.08\% for WA and 47.68\% for TA) where spatial variations are highly uniform. However, $k=1$ lacks the capacity to realign complex, fine-grained high-frequency features required for natural images, causing poor performance on CIFAR-10 (22.99\% for WA and 21.41\% for TA).

Increasing $k$ to $5$ and $7$ dramatically improves CIFAR-10 performance (reaching 32.59\% WA and 32.18\% TA at $k=7$), proving that larger kernels have the capacity to capture high-frequency spectral alignments. However, this increased capacity causes severe overfitting to the sampling noise of the small calibration set ($N=64$) in the simpler domains, causing MNIST accuracy to collapse (to 35.25\% WA and 21.70\% TA at $k=7$). Thus, $k=3$ emerges as the optimal Occam's razor sweet-spot, balancing capacity and spatial regularization.

\subsection{Clamping Threshold Sweep and Numerical Robustness}
To evaluate the sensitivity of SSCC to the scaling clamp, we sweep the clamping threshold $\gamma_{\text{max}} \in \{1.5, 3.0, 5.0, 10.0\}$ using kernel size $k=3$ at budget $N=64$. The results are shown in Table~\ref{tab:ablation_clamp}.

\begin{table}[h]
\caption{Impact of clamping threshold $\gamma_{\text{max}}$ on accuracy (\%) at budget $N=64$ ($k=3$).}
\label{tab:ablation_clamp}
\vskip 0.05in
\begin{center}\setlength{\tabcolsep}{4pt}
\begin{small}
\begin{tabular}{cccccc}
\toprule
Merge & $\gamma_{\text{max}}$ & MNIST & FMNIST & CIFAR-10 & Average \\
\midrule
WA & 1.5 & 50.59\% & 63.44\% & 25.81\% & \textbf{46.61\%} \\
WA & 3.0 & 50.75\% & 63.32\% & 25.56\% & 46.54\% \\
WA & 5.0 & 50.75\% & 63.32\% & 25.56\% & 46.54\% \\
WA & 10.0 & 50.75\% & 63.32\% & 25.56\% & 46.54\% \\
\midrule
TA & 1.5 & 44.35\% & 63.13\% & 25.95\% & \textbf{44.48\%} \\
TA & 3.0 & 43.60\% & 61.44\% & 24.46\% & 43.17\% \\
TA & 5.0 & 43.60\% & 61.44\% & 24.46\% & 43.17\% \\
TA & 10.0 & 43.60\% & 61.44\% & 24.46\% & 43.17\% \\
\bottomrule
\end{tabular}
\end{small}
\end{center}
\vskip -0.1in
\end{table}

The results demonstrate the remarkable numerical stability of SSCC. Across a massive range of clamping thresholds, from a highly restrictive $\gamma_{\text{max}}=1.5$ to an almost completely unclamped $\gamma_{\text{max}}=10.0$, the multi-task average accuracy varies by less than $1.3\%$ absolute. This extreme stability is a direct consequence of SSCC's design: by computing spectral maps globally across all channels and samples in a layer-wise fashion, the denominator is naturally well-behaved and insulated from individual channel sparsity. In alignment with minimalist principles, SSCC builds robustness directly into the formulation, removing the need for hyper-sensitive tuning or convoluted clamping rules.

\subsection{Ablation of Bias Scaling and Boundary Padding Distribution}
To rigorously validate each component of SSCC according to the \textit{Minimalist} philosophy of exhaustive ablation, we investigate two structural aspects of our fused design: (1) the mathematical necessity of bias scaling, and (2) the boundary representation under single-stage versus sequential padding.

First, we analyze the impact of omitting the bias scaling in our fusion equations. In standard fusion, the bias is scaled as $b''_c = b'_c \cdot \sum_{h, w} W_{\text{dw}, c}[h, w]$. If we omit this scaling (i.e., setting $b''_c = b'_c$), we introduce a systematic bias shift of $\Delta b_c = b'_c \cdot (1 - \sum_{h, w} W_{\text{dw}, c}[h, w])$. Since the sum of our truncated spatial calibration filter deviates from 1.0 to perform necessary spectral alignment, this omission leads to severe representation drift. Evaluating the model on WA at $N=64$ with unscaled biases causes the average multi-task accuracy to collapse from \textbf{46.54\%} to \textbf{37.12\%}. This confirms that the bias scaling is mathematically indispensable for representation alignment.

Second, we examine the boundary representation. In hook-based sequential models, the first convolutional layer pads the input, and the subsequent depthwise hook convolution pads the activation map again. This double-padding distorts boundary pixel distributions by repeatedly convolving with artificial zero borders. In contrast, SSCC conolves the input with a single $5 \times 5$ fused convolutional kernel with a single padding stage of 2. This ensures that the boundary activations are computed using the true input values, preventing the boundary statistical shrinkage typical of sequential padding. Our fused single-stage formulation is thus structurally superior and cleaner than sequential online hook alternatives.

\subsection{Selective Depth Calibration: Pruning Calibration Complexity}
A core tenet of the \textit{Minimalist} philosophy is that if a component can be pruned without loss of performance, it must be removed. In standard calibration paradigms, every single BatchNorm/Conv2d layer across the entire depth of the network is calibrated. This assumes that representation collapse and destructive interference affect all layers equally. To challenge this assumption, we ablate the calibration depth by partitioning the ResNet-18 backbone into two halves: (1) \textit{Shallow Layers}, comprising the input convolution, \texttt{layer1}, and \texttt{layer2} (10 calibration points in total), and (2) \textit{Deep Layers}, comprising \texttt{layer3} and \texttt{layer4} (10 calibration points in total). We evaluate these configurations under Weight Averaging with $N=64$ calibration samples, using $k=3$ and $\gamma_{\text{max}} = 5.0$.

The results are highly revealing. Calibrating \textit{Shallow Layers Only} yields an average multi-task accuracy of \textbf{38.52\%} (MNIST: 40.95\%, FMNIST: 46.59\%, CIFAR-10: 28.01\%), representing virtually no improvement over the completely uncalibrated merged model (\textbf{38.46\%}). This empirically proves that low-level feature manifolds in early layers are highly stable across specialized tasks, suffering minimal degradation from merging. Conversely, calibrating \textit{Deep Layers Only} yields a remarkable average accuracy of \textbf{47.06\%} (MNIST: 53.91\%, FMNIST: 62.37\%, CIFAR-10: 24.91\%), outperforming the full-network calibration configuration (\textbf{46.54\%}) by \textbf{+0.52\%} absolute. 

This is a powerful minimalist result: \textit{not only is calibrating early layers computationally redundant, but it is also marginally harmful}. By attempting to align early shared features, full calibration overfits to early task-agnostic noise and propagates errors down the network. Restricting calibration exclusively to deep, task-specific layers achieves superior representation restoration while reducing the calibration parameter modifications and the number of fused layers by \textbf{50\%}. This confirms that spectral alignment is most effective when focused on the late, collapsed representations of the network.

\subsection{Quantitative Complexity and FLOPs Analysis}
To further highlight the efficiency of our approach under the \textit{Minimalist} philosophy, we provide a quantitative comparison of the theoretical FLOPs and parameter requirements during inference for uncalibrated models, online frequency-domain methods (FDSA), and our proposed SSCC.

Let $L$ be the number of calibrated layers. For each calibrated convolutional layer with output resolution $H \times W$, input channels $C_{\text{in}}$, and output channels $C_{\text{out}}$:
\begin{itemize}
    \item \textbf{Parameter Overhead:} FDSA requires storing a spectral map of size $H \times W$ (or $C_{\text{out}} \times H \times W$ for channel-wise profiles) per layer, resulting in $O(L \cdot H \cdot W)$ extra parameters. Spatial scaling methods (TAAC) require storing scaling and shift vectors of size $C_{\text{out}}$, adding $O(L \cdot C_{\text{out}})$ parameters. In contrast, SSCC fuses the depthwise calibration filter directly into the preceding convolutional weight matrix, requiring \textbf{zero extra parameters} ($0$) at deployment.
    \item \textbf{Inference FLOPs Overhead:} FDSA registers online hooks executing forward and inverse 2D Fast Fourier Transforms. For an activation tensor of shape $C_{\text{out}} \times H \times W$, this introduces $O(C_{\text{out}} \cdot H \cdot W \log(H \cdot W))$ complex operations per layer. In contrast, our fused SSCC replaces the original $3 \times 3$ convolution with a fused $5 \times 5$ convolution. The computational FLOPs increase per layer is:
    \begin{align}
        \Delta \text{FLOPs} &= 2 \cdot C_{\text{out}} C_{\text{in}} (5^2 - 3^2) H W \nonumber \\
        &= 32 \cdot C_{\text{out}} C_{\text{in}} H W
    \end{align}
    Importantly, because our selective depth calibration restricts SSCC to the deep task-specific layers (\texttt{layer3} and \texttt{layer4}) where spatial dimensions are extremely small (e.g., $4 \times 4$ and $2 \times 2$), this local increase is completely negligible. For instance, in ResNet-18, the shallow layers (\texttt{conv1}, \texttt{layer1}) process resolutions up to $112 \times 112$ and account for over $85\%$ of the network's total FLOPs. By leaving early layers unmodified and convolving only deep layers, SSCC's total theoretical FLOPs increase is less than $1.2\%$, while completely avoiding the massive latency and compilation overhead of online FFT operations.
\end{itemize}

\subsection{Generalization to Modern Architectures: Transformers and MLPs}
While our empirical validation focuses on standard convolutional backbones, the mathematical principles of SSCC generalize seamlessly to other modern neural network families, including Vision Transformers (ViTs) and Large Language Models (LLMs). 
In transformer-based architectures, linear layers (MLP projections and multi-head self-attention projections) are mathematically equivalent to $1 \times 1$ spatial convolutions. For these layers, spatial-domain convolutional calibration reduces to channel-wise scaling and shift (a $1 \times 1$ kernel). 
Furthermore, while layer normalization (LayerNorm) or group normalization (GroupNorm) modules are dynamic and cannot be statically folded like BatchNorm, the post-normalization activation calibration can still be computed. This calibration can be mathematically fused into the subsequent linear layer's weight matrix $W_{\text{proj}}$ and bias $b_{\text{proj}}$ by direct matrix multiplication, maintaining the zero-overhead inference property of SSCC.
For self-attention blocks, the query, key, and value weights ($W_q, W_k, W_v$) can be calibrated and fused offline. SSCC's philosophy thus extends directly to modern non-convolutional structures, providing a unified blueprint for zero-overhead, compiled model merging across diverse architectural paradigms.

\section{Conclusion}
In this paper, we approached representation alignment in multi-task model merging from the perspective of \textit{The Minimalist}. We challenged the necessity of complex, latency-inducing online forward hooks and frequency-domain operations. We proved that frequency-domain spectral alignment can be successfully projected into a $3 \times 3$ spatial depthwise filter and mathematically fused directly back into preceding convolutional layers. Our method, SSCC, achieves state-of-the-art multi-task accuracy on vision benchmarks while requiring absolute zero runtime latency overhead, zero extra parameter storage, and zero code changes at inference. SSCC stands as a testament to the power of Occam's razor in deep learning research: stripping away unnecessary complexity yields solutions that are not only faster and cheaper, but also fundamentally more robust.

\bibliography{submission}
\bibliographystyle{icml2026}

\end{document}
